\chapter{آليات التفاعل: الأسهم المنحنية}\label{ch:g12:curly-arrows}

\omterm{def:g8:balanced-equations:equation}{المعادلة الموزونة} كالصورتين الأولى والأخيرة من فيلم: فهي
تُظهر \omterm{def:g7:chemical-reactions:reactant}{المتفاعلات} قبلُ \omterm{def:g7:chemical-reactions:reactant}{والنواتج} بعدُ، ولا شيء فيما
بينهما. غير أن \omterm{def:g7:atoms-and-molecules:molecule}{الجزيئات} لا تتبادل \omterm{def:g7:atoms-and-molecules:atom}{الذرات} ببساطة. فالروابط تنكسر وتتكوّن واحدةً
بعد أخرى، بانتقال أزواج \omterm{def:g9:inside-the-atom:nucleus}{الإلكترونات} من موضع إلى آخر، مرورًا بأنواع
كيميائية قصيرة العمر لا تظهر أبدًا في المعادلة. ويرسم الكيميائيون
هذه الحركات بأسهم منحنية. وبقواعد قليلة، تفسّر هذه الأسهم
لماذا يحدث التفاعل حيث يحدث، وتتنبأ بتفاعلات
لم تُشاهد من قبل.

\begin{recall}
\omterm{def:g11:polarity-and-cohesion:electronegativity}{الكهرسلبية} والروابط القطبية والشحنات الجزئية
(\cref{ch:g11:polarity-and-cohesion}). صيغ لويس: الأزواج الرابطة
والأزواج الحرة (\cref{ch:g10:lewis-and-shape}). المجموعات الوظيفية
والعائلات (\cref{ch:g11:functional-groups}).
\end{recall}

\section{المواقع الغنية بالإلكترونات والمواقع الفقيرة إليها}

\begin{definition}[المواقع المانحة للإلكترونات والمستقبلة لها]\label{def:g12:curly-arrows:site}
\emph{الموقع المانح للإلكترونات}\index{موقع مانح للإلكترونات} في \omterm{def:g7:atoms-and-molecules:molecule}{جزيء} أو
\omterm{def:g9:ions:ion}{أيون} موضع غني \omterm{def:g9:inside-the-atom:nucleus}{بالإلكترونات}، قادر على إعطاء زوج لتكوين رابطة
جديدة: \omterm{def:g10:lewis-and-shape:lone-pair}{زوج حر}، أو \omterm{def:g10:lewis-and-shape:multiple-bond}{رابطة ثنائية}، أو \omterm{def:g7:atoms-and-molecules:atom}{ذرة} تحمل شحنة سالبة أو
\omterm{def:g11:polarity-and-cohesion:polar-bond}{شحنة جزئية} $\delta^-$. و\emph{الموقع المستقبِل
للإلكترونات}\index{موقع مستقبل للإلكترونات} موضع فقير إلى \omterm{def:g9:inside-the-atom:nucleus}{الإلكترونات}، قادر على
تلقي زوج: \omterm{def:g7:atoms-and-molecules:atom}{ذرة} تحمل شحنة موجبة أو \omterm{def:g11:polarity-and-cohesion:polar-bond}{شحنة جزئية}
$\delta^+$.
\end{definition}

\begin{example}[المواقع في جزيئين]\label{ex:g12:curly-arrows:sites}
في البرومو إيثان، \ce{CH3-CH2-Br}، البروم أكثر \omterm{def:g11:polarity-and-cohesion:electronegativity}{كهرسلبية} من
الكربون: فالكربون المرتبط به يحمل $\delta^+$ (موقع مستقبِل)،
والبروم $\delta^-$ وثلاثة أزواج حرة. وفي البروبانون، يحمل أكسجين
المجموعة \ce{C=O} شحنة $\delta^-$ وزوجين حرين (موقعان مانحان)، ويحمل
كربونها $\delta^+$ (موقع مستقبِل). \omterm{def:g9:ions:ion}{وأيون} الهيدروكسيد \ce{OH-}، بشحنته
السالبة وأزواجه الحرة الثلاثة، مانح قوي.
\end{example}

\begin{omfigure}
\setchemfig{atom sep=2.2em}
\begin{tabular}{cc}
\chemfig{H_3C-C(-[2]H)(-[6]H)-\charge{90:2pt=\:,0=\:,270:2pt=\:}{Br}} &
\chemfig{H_3C-C(-[6]CH_3)=[2]\charge{0=\:,180=\:}{O}} \\[2pt]
\small $\delta^+$ على الكربون المرتبط بالبروم \ce{Br}، و$\delta^-$ على \ce{Br} &
\small $\delta^+$ على كربون \ce{C=O}، و$\delta^-$ على \ce{O} \\
\small البرومو إيثان & \small البروبانون
\end{tabular}

{\small \omterm{def:g7:atoms-and-molecules:atom}{ذرات} كربون فقيرة إلى \omterm{def:g9:inside-the-atom:nucleus}{الإلكترونات} ($\delta^+$، مواقع مستقبِلة) مجاورة
\omterm{def:g7:atoms-and-molecules:atom}{لذرات} \omterm{def:g11:polarity-and-cohesion:electronegativity}{كهرسلبية} غنية \omterm{def:g9:inside-the-atom:nucleus}{بالإلكترونات} ($\delta^-$، ذات أزواج حرة: مواقع
مانحة).}
\end{omfigure}

\section{الأسهم المنحنية}

\begin{definition}[السهم المنحني]\label{def:g12:curly-arrows:curly-arrow}
\emph{السهم المنحني}\index{سهم منحن} يبيّن حركة زوج من
\omterm{def:g9:inside-the-atom:nucleus}{الإلكترونات} خلال خطوة من تفاعل. ويبدأ من الزوج الذي
يتحرك، زوجًا حرًا كان أو رابطة، وينتهي حيث يذهب الزوج: على \omterm{def:g7:atoms-and-molecules:atom}{ذرة}،
حيث يكوّن رابطة جديدة، أو على \omterm{def:g7:atoms-and-molecules:atom}{ذرة} تأخذ زوج رابطة
تنكسر.
\end{definition}

\begin{method}[رسم الأسهم المنحنية]\label{met:g12:curly-arrows:drawing}
\begin{enumerate}
  \item جد الموقع المانح والموقع المستقبِل.
  \item ارسم سهمًا من الزوج المانح (\omterm{def:g10:lewis-and-shape:lone-pair}{زوج حر} أو رابطة) إلى
        \omterm{def:g7:atoms-and-molecules:atom}{الذرة} المستقبِلة: فتتكوّن هناك رابطة جديدة.
  \item إذا صار \omterm{def:g7:atoms-and-molecules:atom}{للذرة} المستقبِلة عندئذ \omterm{def:g9:inside-the-atom:nucleus}{إلكترونات} أكثر مما ينبغي (أكثر
        من ثمانية، أو أكثر من اثنين للهيدروجين)، فارسم سهمًا ثانيًا
        يكسر إحدى روابطها، ويذهب الزوج إلى \omterm{def:g7:atoms-and-molecules:atom}{الذرة} الأكثر
        \omterm{def:g11:polarity-and-cohesion:electronegativity}{كهرسلبية}.
  \item تحقق من الشحنات بعد الخطوة: \omterm{def:g7:atoms-and-molecules:atom}{الذرة} التي أعطت زوجًا حرًا
        تزيد شحنتها بواحد، \omterm{def:g7:atoms-and-molecules:atom}{والذرة} التي تلقت زوجًا بوصفه
        زوجًا حرًا تنقص شحنتها بواحد؛ والشحنة الكلية لا تتغير.
\end{enumerate}
\end{method}

\begin{proposition}[الشحنة محفوظة في كل خطوة]\label{prop:g12:curly-arrows:charge}
الشحنة الكهربائية الكلية للأنواع الكيميائية هي نفسها قبل كل خطوة
من الآلية وبعدها.
\end{proposition}

\begin{proof}
لا يحرّك \omterm{def:g12:curly-arrows:curly-arrow}{السهم المنحني} إلا \omterm{def:g9:inside-the-atom:nucleus}{إلكترونات} موجودة أصلًا؛ فلا \omterm{def:g9:inside-the-atom:nucleus}{إلكترون}
يُخلق أو يُفنى، والنوى لا تتغير.
\end{proof}

\begin{example}[سهم أول]\label{ex:g12:curly-arrows:first}
يلتقي \omterm{def:g9:ions:ion}{أيون} هيدروجين \ce{H+} \omterm{def:g9:ions:ion}{أيون} هيدروكسيد \ce{OH-}. فيكوّن \omterm{def:g10:lewis-and-shape:lone-pair}{زوج حر} من
الأكسجين رابطة مع \omterm{def:g9:ions:ion}{أيون} الهيدروجين، الذي لا \omterm{def:g9:inside-the-atom:nucleus}{إلكترونات} له:
\ce{H+ + OH- -> H2O}. سهم واحد، من \omterm{def:g10:lewis-and-shape:lone-pair}{الزوج الحر} للأكسجين إلى
\ce{H+}؛ والشحنتان $+1$ و$-1$ تصيران 0 في \omterm{def:g7:atoms-and-molecules:molecule}{جزيء} الماء.
\end{example}

\begin{omfigure}
\setchemfig{atom sep=2.1em}
\babelsublr{\chemfig{@{hp}H^{+}}\qquad\raisebox{0.2ex}{$+$}\qquad
\chemfig{H-@{ox}\charge{90=\:,0=\:,270=\:}{O}^{-}}
\qquad\raisebox{0.2ex}{$\longrightarrow$}\qquad
\chemfig{H-\charge{90=\:,270=\:}{O}-H}}
\chemmove{
  \draw[->, thick, omProp, shorten <=4pt, shorten >=3pt]
    ($(ox)+(0,0.25)$) .. controls +(110:9mm) and +(60:9mm) .. (hp);}

{\small أبسط خطوة: \omterm{def:g10:lewis-and-shape:lone-pair}{زوج حر} من \omterm{def:g9:ions:ion}{أيون} الهيدروكسيد يتحرك ليكوّن
الرابطة الجديدة \ce{O-H} مع \omterm{def:g9:ions:ion}{أيون} الهيدروجين.}
\end{omfigure}

\section{الخطوات الأولية والوسائط}

\begin{definition}[الآلية، الخطوة الأولية، الوسيط]\label{def:g12:curly-arrows:mechanism}
\emph{آلية التفاعل}\index{آلية التفاعل} هي
سلسلة الخطوات التي تصير بها \omterm{def:g7:chemical-reactions:reactant}{المتفاعلات} نواتج. وكل
\emph{خطوة أولية}\index{خطوة أولية} حدث واحد تتحرك فيه
بضعة أزواج من \omterm{def:g9:inside-the-atom:nucleus}{الإلكترونات} دفعة واحدة. \omterm{def:g6:pure-substances-and-mixtures:species}{والنوع الكيميائي} الذي يتكوّن في خطوة
ويُستهلك في خطوة لاحقة هو \emph{وسيط التفاعل}\index{وسيط التفاعل}:
فهو لا يظهر في المعادلة الإجمالية.
\end{definition}


\begin{example}[وسيط كربوكاتيوني]\label{ex:g12:curly-arrows:carbocation}
حين يُضاف بروميد الهيدروجين إلى الإيثين، يجري التفاعل في خطوتين. ففي
الأولى، تعطي \omterm{def:g10:lewis-and-shape:multiple-bond}{الرابطة الثنائية} للإيثين زوجًا لهيدروجين
\ce{HBr}، ويذهب زوج \ce{H-Br} إلى البروم، الذي يغادر
أيونًا \ce{Br-}؛ والكربون الذي فقد حصته من \omterm{def:g10:lewis-and-shape:multiple-bond}{الرابطة الثنائية} يبقى
بستة \omterm{def:g9:inside-the-atom:nucleus}{إلكترونات} فقط وشحنة موجبة: وهو كربوكاتيون،
\ce{CH3-CH2+}. وفي الخطوة الثانية، يكوّن \omterm{def:g10:lewis-and-shape:lone-pair}{زوج حر} من \ce{Br-} رابطة
مع ذلك الكربون. فالكربوكاتيون وسيط.
\end{example}

\section{ثلاث فئات من التفاعلات}

\begin{definition}[الاستبدال، الإضافة، الحذف]\label{def:g12:curly-arrows:categories}
في \emph{الاستبدال}\index{استبدال}، تحل محلَّ \omterm{def:g7:atoms-and-molecules:atom}{ذرة} أو مجموعة في
\omterm{def:g7:atoms-and-molecules:molecule}{جزيء} ذرةٌ أو مجموعة أخرى. وفي \emph{الإضافة}\index{إضافة (تفاعل)}،
يتحد جزيئان في \omterm{def:g7:atoms-and-molecules:molecule}{جزيء} واحد، فتصير رابطة متعددة رابطة أحادية. وفي
\emph{الحذف}\index{حذف (تفاعل)}، يُنتزع \omterm{def:g7:atoms-and-molecules:molecule}{جزيء} صغير (غالبًا الماء)
من \omterm{def:g7:atoms-and-molecules:molecule}{جزيء}، فتصير رابطة أحادية رابطة متعددة.
\end{definition}

\begin{omfigure}
\setchemfig{atom sep=2em}
\small
\textbf{\omterm{def:g12:curly-arrows:categories}{الاستبدال}} (خطوة واحدة): \ce{CH3-CH2-Br + OH- -> CH3-CH2-OH + Br-}

\medskip
\babelsublr{\chemfig{H-@{o1}\charge{90=\:,0=\:,270=\:}{O}^{-}}\qquad
\chemfig{H_3C-@{c1}C(-[2]H)(-[6]H)-[@{b1}]@{br1}\charge{90=\:,0=\:,270=\:}{Br}}
\qquad$\longrightarrow$\qquad
\chemfig{H_3C-C(-[2]H)(-[6]H)-O-H}\quad$+$\quad\ce{Br-}}
\chemmove{
  \draw[->, thick, omProp, shorten <=4pt, shorten >=3pt]
    ($(o1)+(0.15,0.2)$) .. controls +(60:9mm) and +(120:9mm) .. (c1);
  \draw[->, thick, omDef, shorten <=2pt, shorten >=4pt]
    (b1) .. controls +(-90:7mm) and +(-90:7mm) .. (br1);}

\vspace{6mm}
\textbf{الإضافة} (خطوتان): \ce{CH2=CH2 + HBr -> CH3-CH2-Br}

\vspace{9mm}
\babelsublr{\chemfig{H_2C=[@{pi}]CH_2}\quad$+$\quad
\chemfig{@{h2}H-[@{b2}]@{br2}Br}
\qquad$\longrightarrow$\qquad
\chemfig{H_3C-@{cp}\charge{90:2pt=$\scriptstyle +$}{C}H_2}\quad$+$\quad
\chemfig{@{brm}\charge{90=\:,180=\:,270=\:,0=\:}{Br}^{-}}
\qquad$\longrightarrow$\qquad \ce{CH3-CH2-Br}}
\chemmove{
  \draw[->, thick, omProp, shorten <=2pt, shorten >=3pt]
    (pi) .. controls +(90:10mm) and +(110:9mm) .. (h2);
  \draw[->, thick, omDef, shorten <=2pt, shorten >=4pt]
    (b2) .. controls +(-90:7mm) and +(-90:7mm) .. (br2);
  \draw[->, thick, omProp, shorten <=4pt, shorten >=3pt]
    ($(brm)+(0,-0.25)$) .. controls +(-120:8mm) and +(-60:8mm) .. (cp);}

\vspace{6mm}
\textbf{\omterm{def:g12:curly-arrows:categories}{الحذف}} (نزع الماء من الإيثانول في وسط حمضي، ثلاث خطوات):
\ce{CH3-CH2-OH -> CH2=CH2 + H2O}

\medskip
\begin{tabular}{@{}l@{}}
1.\ يأخذ الأكسجين \omterm{def:g9:ions:ion}{أيون} هيدروجين: \ce{CH3-CH2-OH + H+ -> CH3-CH2-OH2+}؛\\
2.\ يغادر زوج \ce{C-O} مع الماء: \ce{CH3-CH2-OH2+ -> CH3-CH2+ + H2O}؛\\
3.\ يصير زوج \ce{C-H} مجاور الرابطةَ الثنائية، ويغادر
    \ce{H+}:\\
\phantom{3.\ }\ce{CH3-CH2+ -> CH2=CH2 + H+}.
\end{tabular}

\medskip
{\small ثلاث آليات. الأسهم البرتقالية: زوج مانح يكوّن رابطة جديدة.
الأسهم الزرقاء: رابطة تنكسر، ويذهب زوجها إلى \omterm{def:g7:atoms-and-molecules:atom}{الذرة} الأكثر
\omterm{def:g11:polarity-and-cohesion:electronegativity}{كهرسلبية}. وفي \omterm{def:g12:curly-arrows:categories}{الحذف} يُعاد \omterm{def:g9:ions:ion}{أيون} الهيدروجين في النهاية.}
\end{omfigure}


\begin{remark}[تغيير الهيكل أو المجموعة]\label{rem:g12:curly-arrows:skeleton}
بعض التفاعلات لا تغيّر إلا \omterm{def:g11:functional-groups:functional-group}{المجموعة الوظيفية} وتحفظ الهيكل
الكربوني: \omterm{def:g12:curly-arrows:categories}{فاستبدال} البرومو إيثان يحوّل ألكانًا مهلجنًا إلى
\omterm{def:g11:functional-groups:alcohol}{كحول}، بذرتي كربون قبله وذرتي كربون بعده. وبعضها الآخر يبني الهيكل أو
يقطعه: فتكوين رابطة كربون--كربون جديدة يطيل السلسلة، وهو الخطوة
الحاسمة في صنع \omterm{def:g7:atoms-and-molecules:molecule}{جزيء} كبير من \omterm{def:g7:atoms-and-molecules:molecule}{جزيئات} صغيرة. وتخطيط التصنيع
يعني اختيار تفاعلات من النوعين، بالترتيب الصحيح.
\end{remark}

\section{تمارين}

\begin{exercise}[$\star$]\label{exo:g12:curly-arrows:1}
سمِّ في كل \omterm{def:g6:pure-substances-and-mixtures:species}{نوع كيميائي} موقعًا مانحًا، وموقعًا مستقبِلًا إن
وُجد: \ce{OH-}؛ \ce{H+}؛ \ce{CH2=CH2}؛ \ce{CH3-Cl}؛ \ce{H2O}.
\end{exercise}

\begin{exercise}[$\star$]\label{exo:g12:curly-arrows:2}
\omterm{def:g12:curly-arrows:categories}{استبدال} أم إضافة أم حذف؟
(a) \ce{CH3-CH2-Cl + OH- -> CH3-CH2-OH + Cl-}؛
(b) \ce{CH2=CH2 + H2O -> CH3-CH2-OH}؛
(c) \ce{CH3-CH2-OH -> CH2=CH2 + H2O}؛
(d) \ce{CH2=CH2 + Br2 -> CH2Br-CH2Br}.
\end{exercise}

\begin{exercise}[$\star$]\label{exo:g12:curly-arrows:3}
من أين يبدأ \omterm{def:g12:curly-arrows:curly-arrow}{السهم المنحني}، وأين ينتهي؟ وماذا يعني سهم منحنٍ
من \omterm{def:g10:lewis-and-shape:lone-pair}{زوج حر} \omterm{def:g7:atoms-and-molecules:atom}{لذرة} أكسجين إلى \omterm{def:g7:atoms-and-molecules:atom}{ذرة} كربون؟
\end{exercise}

\begin{exercise}[$\star$]\label{exo:g12:curly-arrows:4}
في التفاعل \ce{H+ + NH3 -> NH4+}، أي النوعين يعطي الزوج؟
ارسم السهم. وما شحنة الناتج؟
\end{exercise}

\begin{exercise}[$\star$]\label{exo:g12:curly-arrows:5}
ما \omterm{def:g12:curly-arrows:mechanism}{وسيط التفاعل}؟ سمِّ وسيط إضافة
\ce{HBr} إلى الإيثين.
\end{exercise}

\begin{exercise}[$\star\star$]\label{exo:g12:curly-arrows:6}
ارسم الأسهم المنحنية \omterm{def:g12:curly-arrows:categories}{للاستبدال}
\ce{CH3-Br + OH- -> CH3-OH + Br-} (خطوة واحدة)، وتحقق من الشحنات.
\end{exercise}

\begin{exercise}[$\star\star$]\label{exo:g12:curly-arrows:7}
في الخطوة الأولى من إضافة \ce{HBr} إلى الإيثين، لماذا يذهب
زوج \ce{H-Br} إلى البروم لا إلى الهيدروجين؟ وما الشحنة التي
يحملها كل \omterm{def:g7:chemical-reactions:reactant}{ناتج} من نواتج الخطوة؟
\end{exercise}

\begin{exercise}[$\star\star$]\label{exo:g12:curly-arrows:8}
في نزع الماء من الإيثانول، اكتب الخطوات الثلاث وقل، لكل منها،
أي نوع هو المانح وأيها المستقبِل. ولماذا لا يُكتب \omterm{def:g9:ions:ion}{أيون}
الهيدروجين في المعادلة الإجمالية؟
\end{exercise}

\begin{exercise}[$\star\star$]\label{exo:g12:curly-arrows:9}
الخطوة \ce{CH3-CH2+ + H2O -> CH3-CH2-OH2+} تلي
إضافة \omterm{def:g9:ions:ion}{أيون} هيدروجين إلى الإيثين في الماء. ارسم السهم. وماذا
يتكوّن بعد فقد \ce{H+} مرة أخرى؟ اكتب المعادلة الإجمالية
وأعطِ فئتها.
\end{exercise}

\begin{exercise}[$\star\star$]\label{exo:g12:curly-arrows:10}
اقرأ آلية \omterm{def:g12:curly-arrows:categories}{استبدال} البرومو إيثان: كم زوجًا
يتحرك؟ وأي رابطة تتكوّن، وأيها تنكسر؟ وهل هناك وسيط؟
\end{exercise}

\begin{exercise}[$\star\star$]\label{exo:g12:curly-arrows:11}
% ledger: en:C, en:Cl, en:Br
يتفاعل الكلورو ميثان مع \omterm{def:g9:ions:ion}{أيون} الهيدروكسيد أبطأ من
البرومو ميثان. مستعينًا بالكهرسلبيات (C 2.55، Cl 3.16، Br 2.96)،
اشرح أي الكربونين أشد $\delta^+$. وهل يفسّر ذلك
السرعتين؟ (فكّر أيضًا في مدى سهولة أخذ \omterm{def:g9:ions:ion}{الأيون} المغادر للزوج.)
\end{exercise}

\begin{exercise}[$\star\star\star$]\label{exo:g12:curly-arrows:12}
يُضاف كلوريد الهيدروجين إلى البروبين، \ce{CH2=CH-CH3}. وفي الخطوة الأولى
يمكن \omterm{def:g9:ions:ion}{لأيون} الهيدروجين أن يرتبط بأي من كربوني \omterm{def:g10:lewis-and-shape:multiple-bond}{الرابطة الثنائية}، فيعطي
كربوكاتيونين ممكنين. ارسم كليهما، والناتجين الممكنين،
وسمّهما. (ويدرس مجلد السنة الأولى أيهما يتكوّن أكثر.)
\end{exercise}

\begin{exercise}[$\star\star\star$]\label{exo:g12:curly-arrows:13}
يرسم طالب الخطوة الأولى من \omterm{def:g12:curly-arrows:categories}{استبدال} البرومو إيثان بسهم
من \omterm{def:g7:atoms-and-molecules:atom}{ذرة} البروم إلى أكسجين \ce{OH-}. اشرح
الخطأين.
\end{exercise}

\begin{exercise}[$\star\star\star$]\label{exo:g12:curly-arrows:14}
في تكوين \omterm{def:g11:functional-groups:carboxylic-acid}{إستر} من \omterm{def:g11:functional-groups:carboxylic-acid}{حمض كربوكسيلي} \omterm{def:g11:functional-groups:alcohol}{وكحول}، يهاجم
أكسجين \omterm{def:g11:functional-groups:alcohol}{الكحول} كربون المجموعة \ce{C=O}. مستعينًا
بالشحنات الجزئية، اشرح لماذا يكون ذلك الكربون موقعًا مستقبِلًا وذلك
الأكسجين موقعًا مانحًا. وأي \omterm{def:g7:atoms-and-molecules:molecule}{جزيء} صغير يُفقد إجمالًا، وإلى أي
فئة ينتمي التفاعل كله؟
\end{exercise}

\begin{exercise}[$\star\star\star$]\label{exo:g12:curly-arrows:15}
يتفاعل الإيثين مع البروم \ce{Br2}، وهو \omterm{def:g11:polarity-and-cohesion:polar-molecule}{جزيء غير قطبي}. وحين يقترب \omterm{def:g7:atoms-and-molecules:molecule}{جزيء} \ce{Br2}
من \omterm{def:g10:lewis-and-shape:multiple-bond}{الرابطة الثنائية}، تُدفع \omterm{def:g9:inside-the-atom:nucleus}{إلكترونات} \ce{Br-Br}
بعيدًا عن \omterm{def:g7:atoms-and-molecules:atom}{ذرة} البروم القريبة. اشرح لماذا يُنشئ ذلك
موقعًا مستقبِلًا، واقترح الخطوة الأولى من الآلية.
\end{exercise}

\section{مسألة: تحضير البرومو إيثان}

\begin{problem}[{مسألة نهاية الأسبوع --- من عشرة غرامات من الإيثين: كم من
البرومو إيثان؟}]
\label{pb:g12:curly-arrows:1}
% ledger: aw:C, aw:H, aw:Br, en:H, en:Br, en:C
البرومو إيثان، وهو \omterm{def:g5:raw-materials-and-recycling:raw-material}{مادة أولية} لتصنيعات كثيرة، يُحضَّر بإضافة
بروميد الهيدروجين إلى الإيثين: \ce{CH2=CH2 + HBr -> CH3-CH2-Br}. ويستعمل كيميائي
\qty{10.0}{g} من الإيثين وفائضًا من بروميد الهيدروجين؛ \omterm{def:g10:synthesis-yield:yield}{والمردود}
\qty{75}{\percent}.

\medskip
\noindent\textbf{الجزء الأول --- المواقع.}
\begin{enumerate}
  \item ارسم صيغتي لويس للإيثين ولبروميد الهيدروجين.
  \item أين الموقع المانح في الإيثين؟ ولماذا؟
  \item احسب فرق \omterm{def:g11:polarity-and-cohesion:electronegativity}{الكهرسلبية} في الرابطة \ce{H-Br}
        (H 2.20، Br 2.96). أي الذرتين تحمل $\delta^+$؟
  \item أي \omterm{def:g7:atoms-and-molecules:atom}{ذرة} في \ce{HBr} هي الموقع المستقبِل؟
\end{enumerate}

\medskip
\noindent\textbf{الجزء الثاني --- الآلية.}
\begin{enumerate}[resume]
  \item ارسم الخطوة الأولى بسهميها المنحنيين.
  \item أعطِ النوعين المتكوّنين وشحنتيهما.
  \item تحقق من انحفاظ الشحنة في الخطوة الأولى.
  \item ارسم الخطوة الثانية بسهمها المنحني.
  \item كم إلكترونًا يحيط بالكربون الموجب في
        الوسيط؟ ولماذا هو شديد التفاعل؟
\end{enumerate}

\medskip
\noindent\textbf{الجزء الثالث --- أي نوع من التفاعلات؟}
\begin{enumerate}[resume]
  \item ما وسيط الآلية؟
  \item لماذا لا يظهر في المعادلة الإجمالية؟
  \item إلى أي فئة ينتمي التفاعل؟
  \item أيغيّر التفاعل الهيكل الكربوني، أم \omterm{def:g11:functional-groups:functional-group}{المجموعة
        الوظيفية} فقط؟
\end{enumerate}

\medskip
\noindent\textbf{الجزء الرابع --- كم من الناتج؟}
\begin{enumerate}[resume]
  \item احسب الكتلتين الموليتين للإيثين وللبرومو إيثان.
  \item احسب كمية الإيثين.
  \item املأ \omterm{def:g10:reaction-progress-table:extent}{جدول التقدم}. أي المتفاعلين محِدّ؟
  \item احسب الكتلة القصوى للبرومو إيثان.
  \item احسب الكتلة المحصّلة \omterm{def:g10:synthesis-yield:yield}{بمردود} \qty{75}{\percent}.
  \item اذكر الجواب النهائي: ما كتلة البرومو إيثان التي يحصل عليها
        الكيميائي؟
\end{enumerate}
\end{problem}
